\section{Closed nonlinear source images and a dual normal decomposition}
\label{sec:r12-nonlinear-dual-source}

Keep exactly (ED1), the unchanged arbitrary original $S$, and the
Banach space $\mathbb D_S$ of the preceding section. Its positive
potential cone is denoted $\mathscr C$ even though its Borel
representatives are nonpositive functions. Write $J$ for (ED9),
$r(u)=e(u)^{1/d}$, and $\|\cdot\|_*$ for the genuine operator
dual norm (ED10). In particular this is not the absolute-variation
distance of R11 (VF11).

\begin{theorem}[A complete curvature-dual metric and closed source image]
\label{thm:r12-closed-full-source-image}
The image $J(\mathscr C)$ is norm closed in $\mathbb D_S^*$.
The formula
\[
                 d_J(u,v)=\|J(u)-J(v)\|_*
 \tag{NI1}
\]
defines a complete metric on $\mathscr C$, inducing precisely the
R11 strong discrepancy topology. For $u,v\in\mathscr C$,
\[
 \begin{split}
 I(u,v)&\leq(r(u)+r(v))d_J(u,v),\\
 d_J(u,v)&\leq\sqrt q\,
     \max(r(u),r(v))^{(q-1)/2}I(u,v)^{1/2}.
 \end{split}
 \tag{NI2}
\]
No common dominating measure and no fixed actual lower barrier
are required. The assertion concerns all $q\geq1$ and the full
original test category on the stated domain.
\end{theorem}

\noindent\textbf{Status: \texttt{NEW\_PROVED}}. R11 (OI1)--(OI3)
proved closedness only inside a fixed dominated interval for its
stronger absolute-variation distance; neither restriction occurs here.

\begin{proof}
Multilinearity gives the exact identity
$I(u,v)=(J(u)-J(v))(u-v)$. The dual inequality and
$\|u-v\|_{\rm ed}\leq r(u)+r(v)$ prove the first estimate.
The second is (ED11). Thus $d_J=0$ gives $I=0$ and equal
relative classes. Symmetry and the triangle inequality come
from the dual norm, not from an asserted triangle inequality
for $I$.

For completeness, mixed Holder also gives the explicit root bound
\[
 I(u,v)\geq\bigl(r(u)-r(v)\bigr)
                   \bigl(r(u)^q-r(v)^q\bigr)
                \geq|r(u)-r(v)|^{q+1}.
\]
The first inequality follows by bounding the two mixed terms in
$e(u)+e(v)-B(u,v^{[q]})-B(v,u^{[q]})$ separately. The second
is the elementary inequality $x^q-y^q\geq(x-y)^q$ for
$x\geq y\geq0$. Hence convergence to a fixed potential in
$I$ automatically bounds the energy roots.

Suppose $J(u_j)$ is norm Cauchy. Its norms are bounded, so
$r(u_j)^q=\|J(u_j)\|_*$ bounds the energy roots. The first
estimate in (NI2) makes $(u_j)$ strongly Cauchy. R11 strong
completeness gives $u\in\mathscr C$ with $I(u_j,u)\to0$.
Its energy roots remain bounded, and the second estimate gives
$J(u_j)\to J(u)$ in norm. Every norm limit of image elements
is consequently in the image, proving closedness and metric
completeness. The same inequalities prove equivalence of the
two convergences: strong convergence bounds the roots by the
frozen energy-root estimate, whereas $d_J$ convergence bounds
them by the exact dual norm identity. These arguments also
identify their Cauchy sequences.
\end{proof}

\begin{corollary}[Source approximation on every original test]
\label{cor:r12-original-source-approximation}
Let $u_j\in\mathscr C$ and suppose their source functionals
$J(u_j)$ converge in $\mathbb D_S^*$ to $\ell$. Then there is
a unique $u\in\mathscr C$ with $J(u)=\ell$, and $u_j\to u$
strongly. If a specified original Radon measure $\mu$ satisfies
\[
                 \ell(\phi)=-\int_V\phi\,d\mu
            \quad\text{for every original compact smooth scalar }\phi,
 \tag{NI3}
\]
then $\widehat\mu_u=\mu$ on all such tests, after every actual
moment refinement. All other strong current and weighted limits
of R11 (IC11)--(IC13) hold as well.
\end{corollary}

\noindent\textbf{Status: \texttt{DIRECT\_COROLLARY}}.

\begin{proof}
Closedness supplies $u$ and the quantitative estimates supply
strong convergence. Every original compact smooth scalar is an
element of $\mathbb D_S$: R11 (HC14)--(HC15) writes it as the
exact difference of two actual bounded finite-mass common-collar
potentials. That construction is valid for all $q$, not only
for its separate Hilbert-space consequence. Thus (ED9) and
(NI3) give equality on the complete original scalar test sheaf.
Original smooth scalar density and Radon regularity identify
the same measures. The strong calculus supplies the other
limits. This is an approximation criterion, not a claim that
every source admits such image approximants.
\end{proof}

\begin{definition}[The strong-continuous energy dual]
\label{def:r12-strong-continuous-dual}
Let $\mathbb D_{S,\mathrm{s}}^*$ consist of those
$\ell\in\mathbb D_S^*$ for which
\[
  u_j,u\in\mathscr C,\ I(u_j,u)\to0
                 \quad\Longrightarrow\quad\ell(u_j)\to\ell(u).
 \tag{NI4}
\]
This is a continuity requirement on the actual strong cone.
It is not automatically imposed on every Banach-dual functional.
\end{definition}

\begin{lemma}[A closed dual class containing every curvature]
\label{lem:r12-strong-dual-closed}
The class $\mathbb D_{S,\mathrm{s}}^*$ is a norm-closed linear
subspace of $\mathbb D_S^*$, and $J(\mathscr C)$ is contained
in it. Consequently the norm-closed linear span of all curvature
functionals is contained in this class.
\end{lemma}

\noindent\textbf{Status: \texttt{DIRECT\_COROLLARY}} of the
proved Banach realization and frozen strong tester estimate.

\begin{proof}
Linearity is immediate. A strongly convergent row has bounded
$r(u_j)$, hence is bounded in $\mathbb D_S$. If $\ell_m\to\ell$
in norm, the error $|(\ell-\ell_m)(u_j-u)|$ is uniformly small
for large $m$. Fixing $m$ then invokes (NI4), proving closure.
For fixed $v\in\mathscr C$, the completed R11 tester estimate
(EV1) bounds $|J(v)(u_j-u)|$ by
$\int|u_j-u|\,d\widehat\mu_v\to0$. Its constants are finite
on the bounded-energy family. This proves the curvature assertion.
\end{proof}

Let
\[
 \mathscr C^\circ
   =\{n\in\mathbb D_S^*:n(v)\leq0\text{ for every }v\in\mathscr C\}.
 \tag{NI5}
\]
This is the polar of the potential cone in the specified dual,
not an identification of any two current-positivity cones.

\begin{theorem}[All-degree variational normal decomposition]
\label{thm:r12-nonlinear-normal-decomposition}
For every $\ell\in\mathbb D_{S,\mathrm{s}}^*$, the functional
\[
                F_\ell(v)=\ell(v)-e(v)/d\qquad(v\in\mathscr C)
 \tag{NI6}
\]
has a unique maximizer $u=\Pi(\ell)$. Every maximizing sequence
converges strongly to $u$. Equivalently, there is a unique pair
$(u,n)$ with
\[
       \ell=J(u)+n,\qquad
       u\in\mathscr C,\quad n\in\mathscr C^\circ,
                           \quad n(u)=0.
 \tag{NI7}
\]
The residual also lies in $\mathbb D_{S,\mathrm{s}}^*$.
Define the nonnegative support gauge
\[
       p(\ell)=\sup_{\substack{v\in\mathscr C\\r(v)\leq1}}\ell(v).
 \tag{NI8}
\]
Then
\[
 \begin{split}
 e(u)&=\ell(u),\qquad r(u)^q=p(\ell)\leq\|\ell\|_*,\\
 \sup F_\ell&=\frac qd\,p(\ell)^{d/q}.
 \end{split}
 \tag{NI9}
\]
If $\ell$ is nonnegative on $\mathscr C$, then
$p(\ell)=\|\ell\|_*$. No assertion that the residual $n$ vanishes
is part of this theorem.
\end{theorem}

\noindent\textbf{Status: \texttt{NEW\_PROVED}} for arbitrary
strong-continuous signed dual data and every remaining degree.
The midpoint existence mechanism is explicitly inherited from
R11; the Banach source class and exact normal decomposition
are the new extension. This is not unrestricted source solvability.

\begin{proof}
Since $|\ell(v)|\leq\|\ell\|_*r(v)$,
$F_\ell(v)\leq\|\ell\|_*r(v)-r(v)^d/d$. Thus its supremum
$L\geq F_\ell(0)=0$ is finite, and every maximizing sequence
has bounded energies. In the midpoint gap the linear term
cancels. The frozen sharp inequality gives
\[
 c_q I(v_j,v_k)
     \leq L-\tfrac12(F_\ell(v_j)+F_\ell(v_k))\longrightarrow0,
 \quad c_q=\frac{1-2^{-q}}{2d}.
\]
Strong completeness provides $u\in\mathscr C$. Its energy
converges, and (NI4) passes the datum term, so $F_\ell(u)=L$.
The same midpoint estimate proves uniqueness and convergence
of every maximizing sequence. Actual classical strong density
also gives the same supremum over the frozen actual class.

The continuous $d$-linear extension $B$ makes $e/d$ a polynomial
on $\mathbb D_S$ whose differential at $u$ is $J(u)$.
Only segments inside $\mathscr C$ are used for maximization.
Their first variation gives
\[
                   (\ell-J(u))(v-u)\leq0\qquad(v\in\mathscr C).
 \tag{NI10}
\]
Conversely convexity of $e$ on the cone makes this sufficient.
Taking the competitors $0,2u$ gives $(\ell-J(u))(u)=0$.
Then (NI10) is precisely the polar condition in (NI7).
Conversely (NI7) implies (NI10). The pair is unique because
the maximizer is unique and $n=\ell-J(u)$. Its strong continuity
follows from the preceding lemma.

Scaling along $tu$ gives $e(u)=\ell(u)$, also when $u=0$.
For each $v\in\mathscr C$, the polar inequality gives
$\ell(v)\leq J(u)(v)\leq r(v)r(u)^q$. Thus
$p(\ell)\leq r(u)^q$. If $u\ne0$, evaluation on $u/r(u)$
gives the reverse inequality; if $u=0$, the polar condition
and the allowed zero tester give $p(\ell)=0$. This proves
(NI9). Finally (ED10) says
$\|\ell\|_* =\max(p(\ell),p(-\ell))$; if $\ell\geq0$ on
the cone, the latter second term is zero.
\end{proof}

\begin{theorem}[A quantitative nonlinear retraction of source functionals]
\label{thm:r12-source-retraction}
Put $\mathcal P(\ell)=J(\Pi(\ell))$ on
$\mathbb D_{S,\mathrm{s}}^*$. This is a retraction onto the norm-closed
image $J(\mathscr C)$, and is positively homogeneous of degree one.
For $\|\ell\|_*,\|m\|_*\leq R$, with $R>0$, one has
\[
 \begin{split}
 I(\Pi(\ell),\Pi(m))&\leq2R^{1/q}\|\ell-m\|_*,\\
 \|\mathcal P(\ell)-\mathcal P(m)\|_*
                         &\leq\sqrt{2qR}\,\|\ell-m\|_*^{1/2}.
 \end{split}
 \tag{NI11}
\]
The first inequality has the more precise right side
$(\ell-m)(\Pi(\ell)-\Pi(m))$. If
$n=\ell-\mathcal P(\ell)$ and $R=\|\ell\|_*>0$, then
\[
 \left(\sqrt{qR+\|n\|_*}-\sqrt{qR}\right)^2
 \leq\operatorname{dist}_*(\ell,J(\mathscr C))\leq\|n\|_*.
 \tag{NI12}
\]
In particular a nonzero residual is a quantitative obstruction to
source-image approximation, not merely a failure of one chosen
maximization algorithm.
\end{theorem}

\noindent\textbf{Status: \texttt{NEW\_PROVED}} for this nonlinear
all-degree dual retraction and the exact image-approximation gap.
It is not asserted to be a Hilbert metric projection.

\begin{proof}
If $\ell=J(a)$, the pair $(a,0)$ satisfies (NI7), so uniqueness
gives the retraction. Homogeneity of $J$ gives
$\Pi(t\ell)=t^{1/q}\Pi(\ell)$ and
$\mathcal P(t\ell)=t\mathcal P(\ell)$ for $t\geq0$;
the case $t=0$ is the zero class.

Write $u=\Pi(\ell)$ and $v=\Pi(m)$. Use $v,u$ as the
competitors in their two inequalities (NI10), and add. It yields
$I(u,v)\leq(\ell-m)(u-v)$. Duality bounds this by
$\|\ell-m\|_*(r(u)+r(v))$; (NI9) bounds the roots by
$R^{1/q}$. The second line of (NI11) now follows from (ED11).
If $R=0$ both projections are zero.

Let $\delta$ be the distance in (NI12). Since $J(0)=0$,
$\delta\leq R$. Choose $\eta_j\in J(\mathscr C)$ with
$\delta_j=\|\ell-\eta_j\|_*\to\delta$. Their norms are
at most $R+\delta_j$. Retraction and (NI11) imply
\[
 \|n\|_*\leq\delta_j+
               \sqrt{2q(R+\delta_j)\delta_j}.
\]
Passing to the limit and using $\delta\leq R$ gives
$\|n\|_*\leq\delta+2\sqrt{qR\delta}$. Solving this
quadratic inequality in $\sqrt\delta$ gives the lower bound.
The upper bound uses the particular image element $\mathcal P(\ell)$.
\end{proof}

\paragraph{Source equations and scope.}
For an actual measure datum whose literal integration defines $\ell$,
the equation in the completed class holds exactly when the residual
is zero. The decomposition does not assert that an arbitrary positive
Radon datum has that property. The R11 compact-source obstruction
still gives nonzero residuals, and the new higher-degree cap-lift
obstruction supplies further examples. No ambient classical
representative, ordinary pointwise boundary trace, analytic GQ,
or low-real-degree CTO is inferred.
